\section{Taylorreihen}

\subsection{Taylor-Polynom \& Taylor-Reihe}
Das \textbf{Taylor-Polynom} vom Grad \(n\) von \(f\) um \(x_0\):
\[
    p_n(x) = \sum_{k=0}^{n} \frac{f^{(k)}(x_0)}{k!}\,(x-x_0)^k
\]
Die \textbf{Taylor-Reihe} von \(f\) um \(x_0\):
\[
    t_f(x) = \sum_{k=0}^{\infty} \frac{f^{(k)}(x_0)}{k!}\,(x-x_0)^k
\]
Idee: Die Koeffizienten \(a_k = \frac{f^{(k)}(x_0)}{k!}\) werden so gewählt, dass bei \(x_0\) alle Ableitungen bis Grad \(n\) mit \(f\) übereinstimmen.

\subsection{Wichtige Taylor-Reihen (um \texorpdfstring{$x_0=0$}{x0=0})}
\begin{align*}
    e^x     & = \sum_{k=0}^{\infty} \frac{x^k}{k!} = 1 + x + \frac{x^2}{2!} + \frac{x^3}{3!} + \cdots              \\
    \sin(x) & = \sum_{k=0}^{\infty} (-1)^k \frac{x^{2k+1}}{(2k+1)!} = x - \frac{x^3}{3!} + \frac{x^5}{5!} - \cdots \\
    \cos(x) & = \sum_{k=0}^{\infty} (-1)^k \frac{x^{2k}}{(2k)!} = 1 - \frac{x^2}{2!} + \frac{x^4}{4!} - \cdots
\end{align*}

\textbf{Zusammengesetzte Funktionen:} Bekannte Reihe nehmen und Argument ersetzen.\\
Bsp: \(e^{-x^2} = \sum_{k=0}^{\infty} (-1)^k \frac{x^{2k}}{k!}\), \enspace \(\sin(x^3) = \sum_{k=0}^{\infty} (-1)^k \frac{x^{6k+3}}{(2k+1)!}\)

\subsection{Konvergenz}
\(f(x)\) ist durch ihre Taylor-Reihe \textbf{darstellbar}, wenn \(f(x) = t_f(x)\) im \textbf{Konvergenzbereich} gilt, also für alle \(x\), bei denen der Wert der Reihe exakt dem Funktionswert entspricht.

\subsubsection{Potenzreihe (Verallgemeinerung)}
\[
    P(x) = \sum_{k=0}^{\infty} a_k \cdot (x-x_0)^k
\]
Taylor-Reihen sind spezielle Potenzreihen mit \(a_k = \frac{f^{(k)}(x_0)}{k!}\).

\subsubsection{Konvergenzradius}
Für \textit{jede} Potenzreihe existiert ein \(r \geq 0\), sodass:
\begin{itemize}
    \item \(|x-x_0| < r\): Reihe konvergiert
    \item \(|x-x_0| > r\): Reihe divergiert
    \item \(|x-x_0| = r\): keine allgemeine Aussage
\end{itemize}
\(r = \infty\) bedeutet Konvergenz für alle \(x \in \mathbb{R}\).

\subsubsection{Quotientenkriterium}
Berechnung des Konvergenzradius:
\[
   x \in (x_0-r, x_0+r), \quad r = \lim_{k\to\infty} \left|\frac{a_k}{a_{k+1}}\right|
\]

\subsection{Rechnen mit Potenzreihen}
Innerhalb des Konvergenzbereichs wie Polynome behandelbar:
\begin{itemize}
    \item \textbf{Ableitung:} Gliedweise:
          \[P'(x) = \sum_{k=0}^{\infty} (k+1)\,a_{k+1}\,(x-x_0)^k\]
    \item \textbf{Integration:} Gliedweise:
          \[\int P(x)\,dx = \sum_{k=0}^{\infty} \frac{a_k}{k+1}(x-x_0)^{k+1}\]
    \item \textbf{Addition/Subtraktion/Multiplikation} wie bei Polynomen.
\end{itemize}

\subsection{Fehlerabschätzung (Restglied)}
Fehler bei Abbruch nach Grad \(n\):
\[
    |R_n(x)| \leq \left|\frac{(x-x_0)^{n+1}}{(n+1)!}\right| \cdot \max_{z \in [\min(x_0,x),\, \max(x_0,x)]} \left|f^{(n+1)}(z)\right|
\]

\subsection{Anwendungen}

\subsubsection{Numerische Berechnung}
Funktionswerte durch endlich viele Summanden approximieren. Bsp:
\[
    \sin(1.2) \approx 1.2 - \frac{1.2^3}{3!} + \frac{1.2^5}{5!} - \frac{1.2^7}{7!} = 0.9320
\]

\subsubsection{Approximation von Integralen}
Integrand durch Taylor-Polynom ersetzen und integrieren. Bsp:
\[
    \int_0^1 e^{x^2}\,dx \approx \int_0^1 (1+x^2)\,dx = \frac{4}{3}
\]

\subsubsection{Grenzwerte bestimmen}
Funktionen durch Taylor-Reihen ersetzen, kürzen, einsetzen. Bsp:
\[
    \lim_{x\to 0}\frac{1-\cos x}{x^2} = \lim_{x\to 0}\frac{\frac{x^2}{2} - \frac{x^4}{4!} + \cdots}{x^2} = \frac{1}{2}
\]

\subsubsection{Regel von Bernoulli-L'Hôpital}
Für \(\frac{f(x)}{g(x)}\) mit unbestimmtem Ausdruck \(\frac{0}{0}\) oder \(\frac{\infty}{\infty}\) bei \(x \to x_0\):
\[
    \lim_{x\to x_0} \frac{f(x)}{g(x)} = \lim_{x\to x_0} \frac{f'(x)}{g'(x)}
\]
Begründung: Das Verhältnis der Summanden mit kleinstem Exponenten bleibt beim Ableiten erhalten.
